Il metro a nastro con una portata massima di $5\, m$ e una sensibilità di $0.01\,m$ è stato utilizzato per misurare la distanza tra la sorgente di luce e il cellulare (osservatore), quindi per misurare $z$. In questa esperienza ho deciso di usare una torcia elettrica (luce fredda) e la lampadina della mia stanza (luce calda) in modo da visualizzare possibili differenze nelle loro componenti cromatiche. Per reggere la torcia elettrica e il cellulare ho costriuto dei semplici supporti attraverso dei mattoncini Lego in modo da ridurre l'errore associato alla misura di $z$ e fissare in una signola posizione sia la sorgente che "l'osservatore".

L'apparato sperimentale, visualizzabile nell'immagine sottostante, bisogna prepararlo in uno spazio sufficientemente grande in quando se la distanza $z$ è troppo piccola si rischia di non visualizzare il primo ordine di difrazione del secondo reticolo olografico rendendo la misura del suo passo non possibile.

L'apparato va preparato come mostrato nella figura che schematizza l'esperimento. Bisogna porre la sorgente ad una distanza $z$ dal sensore con davanti il reticolo olografico. Particolare attenzione nel porre la sorgente e il sensore alla stessa altezza.

Dopo aver preparato l'apparato sperimentale bisogna fare la prima misura con il metro a nastro. La misura consiste nel misurare la distanza $z$ che separa il sensore (con davanti il reticolo olografico) e la sorgente. Successivamente attraverso il primo reticolo olografico con passo noto si può cominciare il vero e prorpio esperimento. Come detto nella sezione precedente ho deciso di analizzare due diverse sorgenti ma entrambe sono accumunate dal metodo di presa dati. Una volta posizionato il tutto, attraverso il sensore ho preso delle immagini senza reticolo e delle immagini con il reticolo olografico non muovendo nè il sensore nè la sorgente. Le immagini senza reticolo le ho prese con la luce accesa in modo da visualizzare al meglio il supporto della sorgente che prendo come riferimento di misura mentre le immagini con il reticolo olografico le ho prese con tutte le luci della stanza spente e impostando dal cellulare un basso tempo di esposizione in modo tale da catturare la componenti cromatiche al meglio. 

Dopo aver ottenuto tutte le foto con entrambi i reticoli olografici ho spostato le foto sul PC e scaricato il sofrware imagej per analizzare e misurare le distanze tra la sorgente e le proprie componenti cromatiche. Mettendo come misura di riferimento il supporto della sorgente misurato in precedenza e visualizzabile in maniera chiara nelle foto prese con la luce accesa sono riuscito a misurare la distanza tra la sorgente e le componenti cromatiche del primo ordine. Per il primo reticolo olografico misurare la distanza $\Delta x$ serve a ricavare la lunghezza d'onda $\lambda$ utilizzando la formula (\ref{deltax}). Successivamente con lo stesso procedimento utilizzando il software Imagej si misura di nuovo il $\Delta x$ dall'immagine del secondo reticolo olografico con $d$ sconosciuta e attraverso $\lambda$ misurata precedentemente e $\Delta x$ con la formula (\ref{d}) si ottiene il passo.